% ****************************************************************************************************
% classicthesis-config.tex 
% formerly known as loadpackages.sty, classicthesis-ldpkg.sty, and classicthesis-preamble.sty 
% Use it at the beginning of your ClassicThesis.tex, or as a LaTeX Preamble 
% in your ClassicThesis.{tex,lyx} with % ****************************************************************************************************
% classicthesis-config.tex 
% formerly known as loadpackages.sty, classicthesis-ldpkg.sty, and classicthesis-preamble.sty 
% Use it at the beginning of your ClassicThesis.tex, or as a LaTeX Preamble 
% in your ClassicThesis.{tex,lyx} with % ****************************************************************************************************
% classicthesis-config.tex 
% formerly known as loadpackages.sty, classicthesis-ldpkg.sty, and classicthesis-preamble.sty 
% Use it at the beginning of your ClassicThesis.tex, or as a LaTeX Preamble 
% in your ClassicThesis.{tex,lyx} with % ****************************************************************************************************
% classicthesis-config.tex 
% formerly known as loadpackages.sty, classicthesis-ldpkg.sty, and classicthesis-preamble.sty 
% Use it at the beginning of your ClassicThesis.tex, or as a LaTeX Preamble 
% in your ClassicThesis.{tex,lyx} with \input{classicthesis-config}
% ****************************************************************************************************  
% If you like the classicthesis, then I would appreciate a postcard. 
% My address can be found in the file ClassicThesis.pdf. A collection 
% of the postcards I received so far is available online at 
% http://postcards.miede.de
% ****************************************************************************************************


% ****************************************************************************************************
% 0. Set the encoding of your files. UTF-8 is the only sensible encoding nowadays. If you can't read
% äöüßáéçèê∂åëæƒÏ€ then change the encoding setting in your editor, not the line below. If your editor
% does not support utf8 use another editor!
% ****************************************************************************************************
\PassOptionsToPackage{utf8}{inputenc}
	\usepackage{inputenc}

% ****************************************************************************************************
% 1. Configure classicthesis for your needs here, e.g., remove "drafting" below 
% in order to deactivate the time-stamp on the pages
% ****************************************************************************************************
\PassOptionsToPackage{eulerchapternumbers,listings,%
					 pdfspacing,%floatperchapter,%linedheaders,%
					 subfig,beramono,parts}{classicthesis}                                        
% ********************************************************************
% Available options for classicthesis.sty 
% (see ClassicThesis.pdf for more information):
% drafting
% parts nochapters linedheaders
% eulerchapternumbers beramono eulermath pdfspacing minionprospacing
% tocaligned dottedtoc manychapters
% listings floatperchapter subfig
% ********************************************************************


% ****************************************************************************************************
% 2. Personal data and user ad-hoc commands
% ****************************************************************************************************
\newcommand{\myTitle}{Real-World Reinforcement Learning for Vision-Language Robotic Manipulation\xspace}
\newcommand{\mySubtitle}{Deep-learning for Robotics\xspace}
\newcommand{\myDegree}{Bachelor of Engineering Honours and Computer Science\xspace}
\newcommand{\myName}{Lucas Barbosa\xspace}
\newcommand{\myProf}{Will Midgley\xspace}
\newcommand{\myOtherProf}{Put name here\xspace}
\newcommand{\mySupervisor}{Will Midgley\xspace}
\newcommand{\myFaculty}{Faculty of Engineering\xspace}
\newcommand{\myDepartment}{School of Mechanical and Manufacturing Engineering\xspace}
\newcommand{\myUni}{UNSW\xspace}
\newcommand{\myLocation}{Sydney\xspace}
% \newcommand{\myTime}{November 2025\xspace}
% \newcommand{\myVersion}{version 4.2\xspace}

% ********************************************************************
% Setup, finetuning, and useful commands
% ********************************************************************
\newlength{\abcd} % for ab..z string length calculation
\providecommand{\mLyX}{L\kern-.1667em\lower.25em\hbox{Y}\kern-.125emX\@}
\newcommand{\ie}{i.\,e.}
\newcommand{\Ie}{I.\,e.}
\newcommand{\eg}{e.\,g.}
\newcommand{\Eg}{E.\,g.} 
% ****************************************************************************************************


% ****************************************************************************************************
% 3. Loading some handy packages
% ****************************************************************************************************
% ******************************************************************** 
% Packages with options that might require adjustments
% ******************************************************************** 
%\PassOptionsToPackage{ngerman,american}{babel}   % change this to your language(s)
% Spanish languages need extra options in order to work with this template
%\PassOptionsToPackage{spanish,es-lcroman}{babel}
	\usepackage{babel}                  

\usepackage{csquotes}
\PassOptionsToPackage{%
    backend=biber,%
	language=auto,%
	style=numeric-comp,%
    %style=authoryear-comp, % Author 1999, 2010
    %bibstyle=authoryear,dashed=false, % dashed: substitute rep. author with ---
    sorting=nyt, % name, year, title
    maxbibnames=10, % default: 3, et al.
    %backref=true,%
    natbib=true % natbib compatibility mode (\citep and \citet still work)
}{biblatex}
    \usepackage{biblatex}

% Maths
\usepackage{amsmath, amssymb}

% ******************************************************************** 
% General useful packages
% ******************************************************************** 
\PassOptionsToPackage{T1}{fontenc} % T2A for cyrillics
    \usepackage{fontenc}     
\usepackage{textcomp} % fix warning with missing font shapes
\usepackage{scrhack} % fix warnings when using KOMA with listings package          
\usepackage{xspace} % to get the spacing after macros right  
\usepackage{mparhack} % get marginpar right
%\usepackage[latest]{latexrelease} % will be used once available in more distributions (ISSUE #107)
\PassOptionsToPackage{printonlyused,smaller}{acronym} 
    \usepackage{acronym} % nice macros for handling all acronyms in the thesis
    %\renewcommand{\bflabel}[1]{{#1}\hfill} % fix the list of acronyms --> no longer working
    %\renewcommand*{\acsfont}[1]{\textsc{#1}} 
    \renewcommand*{\aclabelfont}[1]{\acsfont{#1}}
% ****************************************************************************************************


% ****************************************************************************************************
% 4. Setup floats: tables, (sub)figures, and captions
% ****************************************************************************************************
\usepackage{tabularx} % better tables
    \setlength{\extrarowheight}{3pt} % increase table row height
\newcommand{\tableheadline}[1]{\multicolumn{1}{c}{\spacedlowsmallcaps{#1}}}
\newcommand{\myfloatalign}{\centering} % to be used with each float for alignment
% \usepackage{caption}
% \usepackage[labelsep=period, labelformat=default,font=footnotesize,labelfont=bf, width=\textwidth,format=plain]{caption}
% Thanks to cgnieder and Claus Lahiri
% http://tex.stackexchange.com/questions/69349/spacedlowsmallcaps-in-caption-label
% [REMOVED DUE TO OTHER PROBLEMS, SEE ISSUE #82]    
%\DeclareCaptionLabelFormat{smallcaps}{\bothIfFirst{#1}{~}\MakeTextLowercase{\textsc{#2}}}
%\captionsetup{font=small,labelformat=smallcaps} % format=hang,
\usepackage[labelsep=period,labelformat=default,font=footnotesize,labelfont=bf,format=plain]{caption}
\usepackage{subfig}
% ****************************************************************************************************


% ****************************************************************************************************
% 5. Setup code listings
% ****************************************************************************************************
\usepackage{listings} 
%\lstset{emph={trueIndex,root},emphstyle=\color{BlueViolet}}%\underbar} % for special keywords
\lstset{language=[LaTeX]Tex,%C++,
    morekeywords={PassOptionsToPackage,selectlanguage},
    keywordstyle=\color{RoyalBlue},%\bfseries,
    basicstyle=\small\ttfamily,
    %identifierstyle=\color{NavyBlue},
    commentstyle=\color{Green}\ttfamily,
    stringstyle=\rmfamily,
    numbers=none,%left,%
    numberstyle=\scriptsize,%\tiny
    stepnumber=5,
    numbersep=8pt,
    showstringspaces=false,
    breaklines=true,
    %frameround=ftff,
    %frame=single,
    belowcaptionskip=.75\baselineskip
    %frame=L
} 
% ****************************************************************************************************             

% ****************************************************************************************************
% 6. PDFLaTeX, hyperreferences and citation backreferences
% ****************************************************************************************************
% ********************************************************************
% Using PDFLaTeX
% ********************************************************************
\usepackage{titlesec} % Let hyperref place anchors with the typeset headings.
\PassOptionsToPackage{pdftex,hyperfootnotes=false,pdfpagelabels}{hyperref}
    \usepackage{hyperref}  % backref linktocpage pagebackref
\pdfcompresslevel=9
\pdfadjustspacing=1 
\PassOptionsToPackage{pdftex}{graphicx}
    \usepackage{graphicx} 
\usepackage[type=CC,modifier=by-nc-nd,version=4.0,imagewidth=25mm]{doclicense}
\renewcommand*{\doclicenseURL}{https://creativecommons.org/licenses/by-nc-nd/4.0/}
\newcommand*{\myLicenseNotice}{Except where otherwise indicated, the original
    text and illustrations in \textit{\myTitle} by \myName{} and contributors
    are licensed under the \doclicenseLongName{} License}
\renewcommand*{\doclicenseLongTextForHyperref}{\myLicenseNotice.}

\newenvironment{spacedparagraphs}{%
    \par
    \setlength{\parindent}{0pt}%
    \setlength{\parskip}{0.75\baselineskip}%
}{\par}

% ********************************************************************
% Hyperreferences
% ********************************************************************
\hypersetup{%
    %draft, % = no hyperlinking at all (useful in b/w printouts)
    colorlinks=true, linktocpage=true, pdfstartpage=3, pdfstartview={FitV null},%
    % uncomment the following line if you want to have black links (e.g., for printing)
    %colorlinks=false, linktocpage=false, pdfstartpage=3, pdfstartview=FitV, pdfborder={0 0 0},%
    breaklinks=true, pdfpagemode=UseNone, pageanchor=true, pdfpagemode=UseOutlines,%
    plainpages=false, bookmarksnumbered, bookmarksopen=true, bookmarksopenlevel=1,%
    hypertexnames=true, pdfhighlight=/O,%nesting=true,%frenchlinks,%
    urlcolor=webbrown, linkcolor=RoyalBlue, citecolor=webgreen, %pagecolor=RoyalBlue,%
    %urlcolor=Black, linkcolor=Black, citecolor=Black, %pagecolor=Black,%
    keeppdfinfo, % Include the author in both PDF properties and XMP metadata.
    pdftitle={\myTitle},%
    pdfauthor={\myName},%
    pdfsubject={},%
    pdfkeywords={},%
    pdfcreator={pdfLaTeX},%
    pdfproducer={LaTeX with hyperref and classicthesis}%
}   

% ********************************************************************
% Setup autoreferences
% ********************************************************************
% There are some issues regarding autorefnames
% http://www.ureader.de/msg/136221647.aspx
% http://www.tex.ac.uk/cgi-bin/texfaq2html?label=latexwords
% you have to redefine the makros for the 
% language you use, e.g., american, ngerman
% (as chosen when loading babel/AtBeginDocument)
% ********************************************************************
\makeatletter
\@ifpackageloaded{babel}%
    {%
       \addto\extrasamerican{%
			\renewcommand*{\figureautorefname}{Figure}%
			\renewcommand*{\tableautorefname}{Table}%
			\renewcommand*{\partautorefname}{Part}%
			\renewcommand*{\chapterautorefname}{Chapter}%
			\renewcommand*{\sectionautorefname}{Section}%
			\renewcommand*{\subsectionautorefname}{Section}%
			\renewcommand*{\subsubsectionautorefname}{Section}%     
                }%
       \addto\extrasngerman{% 
			\renewcommand*{\paragraphautorefname}{Absatz}%
			\renewcommand*{\subparagraphautorefname}{Unterabsatz}%
			\renewcommand*{\footnoteautorefname}{Fu\"snote}%
			\renewcommand*{\FancyVerbLineautorefname}{Zeile}%
			\renewcommand*{\theoremautorefname}{Theorem}%
			\renewcommand*{\appendixautorefname}{Anhang}%
			\renewcommand*{\equationautorefname}{Gleichung}%        
			\renewcommand*{\itemautorefname}{Punkt}%
                }%  
            % Fix to getting autorefs for subfigures right (thanks to Belinda Vogt for changing the definition)
            \providecommand{\subfigureautorefname}{\figureautorefname}%             
    }{\relax}
\makeatother


% ****************************************************************************************************
% 7. Last calls before the bar closes
% ****************************************************************************************************

% ********************************************************************
% Development Stuff
% ********************************************************************
\listfiles
%\PassOptionsToPackage{l2tabu,orthodox,abort}{nag}
%   \usepackage{nag}
%\PassOptionsToPackage{warning, all}{onlyamsmath}
%   \usepackage{onlyamsmath}

% ********************************************************************
% Last, but not least...
% ********************************************************************
\usepackage{microtype} % Improves typography, reduces overflow via character protrusion/expansion
\usepackage[dottedtoc]{classicthesis} 
\usepackage{lmodern}

% ****************************************************************************************************

% ****************************************************************************************************
% 8. Further adjustments (experimental)
% ****************************************************************************************************
% ********************************************************************
% Changing the text area
% ********************************************************************
\linespread{1.05} % a bit more for Palatino
% \areaset[current]{312pt}{761pt} % 686 (factor 2.2) + 33 head + 42 head \the\footskip
% \setlength{\marginparwidth}{7em}%
% \setlength{\marginparsep}{2em}%

% ********************************************************************
% Fix text/equation overflow into margins
% ********************************************************************
\tolerance=1500                    % Allow looser spacing (default 200)
\emergencystretch=1.5em            % Extra stretch to avoid overflow
\hfuzz=1.5pt                       % Suppress warnings for minor overflow
\hbadness=1500   

% ********************************************************************
% Using different fonts
% ********************************************************************
%\usepackage[oldstylenums]{kpfonts} % oldstyle notextcomp
%\usepackage[osf]{libertine}
%\usepackage[light,condensed,math]{iwona}
%\renewcommand{\sfdefault}{iwona}
%\usepackage{lmodern} % <-- no osf support :-(
%\usepackage{cfr-lm} % 
%\usepackage[urw-garamond]{mathdesign} <-- no osf support :-(
%\usepackage[default,osfigures]{opensans} % scale=0.95 
%\usepackage[sfdefault]{FiraSans}
% ****************************************************************************************************

% ****************************************************************************************************
% X. My Custom Stuff
% ****************************************************************************************************
\usepackage{nomencl}
\usepackage{setspace}        % for \onehalfspacing
\usepackage{svg}
\usepackage{epigraph}
\usepackage{float}
\usepackage[section]{placeins} % keep floats within their sections/chapters
\usepackage{booktabs}
\usepackage{import}
\usepackage{array}
\usepackage{ragged2e}
\usepackage{longtable}
\usepackage{multirow}
\setlength{\LTleft}{0pt} % Fix longtable shift on page breaks in ClassicThesis
\setlength{\LTright}{\fill}
\usepackage{tikz}
\usepackage{lipsum}
% \usepackage[scaled=0.85]{beramono} % reload, after lmodern
% \renewcommand{\ttdefault}{fvm}     % make sure \ttfamily = Bera Mono
\usepackage[table]{xcolor}
\usepackage[most]{tcolorbox}
\usepackage[para]{threeparttable}
\usepackage{adjustbox}

\usepackage{xkeyval}     % for \define@key[FigTools]{...}, \setkeys
\usepackage{ifthen}      % for \ifthenelse and \equal
\usepackage{ifoddpage}   % for \checkoddpage and \ifoddpage

% Algorithms
\usepackage{algorithm}
\usepackage{algpseudocode}
% Each algorithm restarts its line numbers; include the float in PDF destinations.
\counterwithin*{ALG@line}{algorithm}

\usepackage{url}

% ********************************************************************
% CLEVEREF SETUP
% ********************************************************************
\usepackage[nameinlink,capitalize]{cleveref}
% Older LaTeX formats lose all but the first letter of acronym labels through
% cleveref's amsmath adapter. Use the complete key for acronym first-use links.
% https://github.com/oetiker/acronym/issues/62
\makeatletter
\IfFormatAtLeastTF{2025-06-01}{}{%
    \patchcmd{\AC@verridelabel}
        {\ifdefined\ltx@label\ltx@label}
        {\ifdefined\cref@label\cref@label}{}{}%
}
\makeatother

% 1. GENERIC SETTINGS
% -------------------
% Capitalize all references (even mid-sentence) to match Engineering standards
% e.g., "as seen in Fig. 1" rather than "as seen in fig. 1"

% 2. SPECIFIC DEFINITIONS
% -----------------------

% Equations: "Eq." and "Eqs."
\crefname{equation}{Eq.}{Eqs.}
\Crefname{equation}{Eq.}{Eqs.}
% Optional: Remove parentheses around eq number (e.g., "Eq. 1" instead of "Eq. (1)")
% If you prefer standard math format "Eq. (1)", comment out the line below.
\creflabelformat{equation}{#2#1#3}

% Figures: "Fig." and "Figs."
\crefname{figure}{Fig.}{Figs.}
\Crefname{figure}{Fig.}{Figs.}

% Tables: "Table" and "Tables" (Usually fully written out)
\crefname{table}{Table}{Tables}
\Crefname{table}{Table}{Tables}

% Code Listings (from listings package): "Listing" and "Listings"
\crefname{lstlisting}{Listing}{Listings}
\Crefname{lstlisting}{Listing}{Listings}

% Algorithms (from algorithm package): "Algorithm" and "Algorithms"
\crefname{algorithm}{Algorithm}{Algorithms}
\Crefname{algorithm}{Algorithm}{Algorithms}

% Sections/Chapters: "Section", "Chapter", "Part"
\crefname{part}{Part}{Parts}
\Crefname{part}{Part}{Parts}
\renewcommand{\thepart}{\Roman{part}}
% A ClassicThesis part ships its title page before \part returns.
\makeatletter
\newcommand{\labeledpart}[2]{%
    \GetTitleString{#1}\let\@currentlabelname\GetTitleStringResult
    \part[#1]{#1\label{#2}}%
}
\makeatother

\crefname{chapter}{Chapter}{Chapters}
\Crefname{chapter}{Chapter}{Chapters}

\crefname{section}{Section}{Sections}
\Crefname{section}{Section}{Sections}

\crefname{subsection}{Section}{Sections} % Map subsections to "Section"
\Crefname{subsection}{Section}{Sections}

\crefname{subsubsection}{Section}{Sections} % Map sub-subsections to "Section"
\Crefname{subsubsection}{Section}{Sections}

% Appendices
\crefname{appendix}{Appendix}{Appendices}
\Crefname{appendix}{Appendix}{Appendices}

% Enumerated Lists (items)
\crefname{enumi}{Item}{Items}
\Crefname{enumi}{Item}{Items}
\crefname{enumii}{Item}{Items}
\Crefname{enumii}{Item}{Items}
\crefname{enumiii}{Item}{Items}
\Crefname{enumiii}{Item}{Items}

% url penalties
\setcounter{biburllcpenalty}{7000}
\setcounter{biburlucpenalty}{8000}

% your column types
\newcolumntype{L}[1]{>{\RaggedRight\arraybackslash}p{#1}}
\newcolumntype{R}[1]{>{\RaggedLeft\arraybackslash}p{#1}}

% tcolorbox style
% Define color themes
\definecolor{lucasredlight}{RGB}{255,230,230}
\definecolor{lucasreddark}{RGB}{180,0,0}

\definecolor{lucasbluelight}{RGB}{225,235,255}
\definecolor{lucasbluedark}{RGB}{0,60,150}

% A key to choose theme = red or blue
\tcbset{
  lucasbox/.style={
    enhanced,
    fonttitle=\bfseries,
    title filled,
    boxrule=1pt,
    toptitle=2mm,
    bottomtitle=1mm
  },
  lucasboxtheme/red/.style={
    colback=lucasredlight,
    colframe=lucasreddark,
    coltitle=white,
    colbacktitle=lucasreddark
  },
  lucasboxtheme/blue/.style={
    colback=lucasbluelight,
    colframe=lucasbluedark,
    coltitle=white,
    colbacktitle=lucasbluedark
  }
}

% Default theme = red
\newenvironment{importantbox}[2][red]
  {\begin{tcolorbox}[lucasbox, lucasboxtheme/#1, title=#2]}
  {\end{tcolorbox}}

\DeclareMathOperator*{\argmax}{arg\,max}

\newcommand{\quotegraffito}[2]{\graffito{#1\\--~\emph{#2}}}

% grey midrule for trial block boundaries
\newcommand{\graymidrule}{%
  \arrayrulecolor{gray!60}%
  \midrule
  \arrayrulecolor{black}%
}

% Method labels (math and text safe)
\newcommand{\mBase}{\text{\textsc{Base}}}   % pretrained / SFT base policy
\newcommand{\mRL}{\text{\textsc{RL}}}       % pure RL from scratch
\newcommand{\mResRL}{\text{\textsc{ResRL}}} % residual RL on Base

% ****************************************************************************************************  
% If you like the classicthesis, then I would appreciate a postcard. 
% My address can be found in the file ClassicThesis.pdf. A collection 
% of the postcards I received so far is available online at 
% http://postcards.miede.de
% ****************************************************************************************************


% ****************************************************************************************************
% 0. Set the encoding of your files. UTF-8 is the only sensible encoding nowadays. If you can't read
% äöüßáéçèê∂åëæƒÏ€ then change the encoding setting in your editor, not the line below. If your editor
% does not support utf8 use another editor!
% ****************************************************************************************************
\PassOptionsToPackage{utf8}{inputenc}
	\usepackage{inputenc}

% ****************************************************************************************************
% 1. Configure classicthesis for your needs here, e.g., remove "drafting" below 
% in order to deactivate the time-stamp on the pages
% ****************************************************************************************************
\PassOptionsToPackage{eulerchapternumbers,listings,%
					 pdfspacing,%floatperchapter,%linedheaders,%
					 subfig,beramono,parts}{classicthesis}                                        
% ********************************************************************
% Available options for classicthesis.sty 
% (see ClassicThesis.pdf for more information):
% drafting
% parts nochapters linedheaders
% eulerchapternumbers beramono eulermath pdfspacing minionprospacing
% tocaligned dottedtoc manychapters
% listings floatperchapter subfig
% ********************************************************************


% ****************************************************************************************************
% 2. Personal data and user ad-hoc commands
% ****************************************************************************************************
\newcommand{\myTitle}{Real-World Reinforcement Learning for Vision-Language Robotic Manipulation\xspace}
\newcommand{\mySubtitle}{Deep-learning for Robotics\xspace}
\newcommand{\myDegree}{Bachelor of Engineering Honours and Computer Science\xspace}
\newcommand{\myName}{Lucas Barbosa\xspace}
\newcommand{\myProf}{Will Midgley\xspace}
\newcommand{\myOtherProf}{Put name here\xspace}
\newcommand{\mySupervisor}{Will Midgley\xspace}
\newcommand{\myFaculty}{Faculty of Engineering\xspace}
\newcommand{\myDepartment}{School of Mechanical and Manufacturing Engineering\xspace}
\newcommand{\myUni}{UNSW\xspace}
\newcommand{\myLocation}{Sydney\xspace}
% \newcommand{\myTime}{November 2025\xspace}
% \newcommand{\myVersion}{version 4.2\xspace}

% ********************************************************************
% Setup, finetuning, and useful commands
% ********************************************************************
\newlength{\abcd} % for ab..z string length calculation
\providecommand{\mLyX}{L\kern-.1667em\lower.25em\hbox{Y}\kern-.125emX\@}
\newcommand{\ie}{i.\,e.}
\newcommand{\Ie}{I.\,e.}
\newcommand{\eg}{e.\,g.}
\newcommand{\Eg}{E.\,g.} 
% ****************************************************************************************************


% ****************************************************************************************************
% 3. Loading some handy packages
% ****************************************************************************************************
% ******************************************************************** 
% Packages with options that might require adjustments
% ******************************************************************** 
%\PassOptionsToPackage{ngerman,american}{babel}   % change this to your language(s)
% Spanish languages need extra options in order to work with this template
%\PassOptionsToPackage{spanish,es-lcroman}{babel}
	\usepackage{babel}                  

\usepackage{csquotes}
\PassOptionsToPackage{%
    backend=biber,%
	language=auto,%
	style=numeric-comp,%
    %style=authoryear-comp, % Author 1999, 2010
    %bibstyle=authoryear,dashed=false, % dashed: substitute rep. author with ---
    sorting=nyt, % name, year, title
    maxbibnames=10, % default: 3, et al.
    %backref=true,%
    natbib=true % natbib compatibility mode (\citep and \citet still work)
}{biblatex}
    \usepackage{biblatex}

% Maths
\usepackage{amsmath, amssymb}

% ******************************************************************** 
% General useful packages
% ******************************************************************** 
\PassOptionsToPackage{T1}{fontenc} % T2A for cyrillics
    \usepackage{fontenc}     
\usepackage{textcomp} % fix warning with missing font shapes
\usepackage{scrhack} % fix warnings when using KOMA with listings package          
\usepackage{xspace} % to get the spacing after macros right  
\usepackage{mparhack} % get marginpar right
%\usepackage[latest]{latexrelease} % will be used once available in more distributions (ISSUE #107)
\PassOptionsToPackage{printonlyused,smaller}{acronym} 
    \usepackage{acronym} % nice macros for handling all acronyms in the thesis
    %\renewcommand{\bflabel}[1]{{#1}\hfill} % fix the list of acronyms --> no longer working
    %\renewcommand*{\acsfont}[1]{\textsc{#1}} 
    \renewcommand*{\aclabelfont}[1]{\acsfont{#1}}
% ****************************************************************************************************


% ****************************************************************************************************
% 4. Setup floats: tables, (sub)figures, and captions
% ****************************************************************************************************
\usepackage{tabularx} % better tables
    \setlength{\extrarowheight}{3pt} % increase table row height
\newcommand{\tableheadline}[1]{\multicolumn{1}{c}{\spacedlowsmallcaps{#1}}}
\newcommand{\myfloatalign}{\centering} % to be used with each float for alignment
% \usepackage{caption}
% \usepackage[labelsep=period, labelformat=default,font=footnotesize,labelfont=bf, width=\textwidth,format=plain]{caption}
% Thanks to cgnieder and Claus Lahiri
% http://tex.stackexchange.com/questions/69349/spacedlowsmallcaps-in-caption-label
% [REMOVED DUE TO OTHER PROBLEMS, SEE ISSUE #82]    
%\DeclareCaptionLabelFormat{smallcaps}{\bothIfFirst{#1}{~}\MakeTextLowercase{\textsc{#2}}}
%\captionsetup{font=small,labelformat=smallcaps} % format=hang,
\usepackage[labelsep=period,labelformat=default,font=footnotesize,labelfont=bf,format=plain]{caption}
\usepackage{subfig}
% ****************************************************************************************************


% ****************************************************************************************************
% 5. Setup code listings
% ****************************************************************************************************
\usepackage{listings} 
%\lstset{emph={trueIndex,root},emphstyle=\color{BlueViolet}}%\underbar} % for special keywords
\lstset{language=[LaTeX]Tex,%C++,
    morekeywords={PassOptionsToPackage,selectlanguage},
    keywordstyle=\color{RoyalBlue},%\bfseries,
    basicstyle=\small\ttfamily,
    %identifierstyle=\color{NavyBlue},
    commentstyle=\color{Green}\ttfamily,
    stringstyle=\rmfamily,
    numbers=none,%left,%
    numberstyle=\scriptsize,%\tiny
    stepnumber=5,
    numbersep=8pt,
    showstringspaces=false,
    breaklines=true,
    %frameround=ftff,
    %frame=single,
    belowcaptionskip=.75\baselineskip
    %frame=L
} 
% ****************************************************************************************************             

% ****************************************************************************************************
% 6. PDFLaTeX, hyperreferences and citation backreferences
% ****************************************************************************************************
% ********************************************************************
% Using PDFLaTeX
% ********************************************************************
\usepackage{titlesec} % Let hyperref place anchors with the typeset headings.
\PassOptionsToPackage{pdftex,hyperfootnotes=false,pdfpagelabels}{hyperref}
    \usepackage{hyperref}  % backref linktocpage pagebackref
\pdfcompresslevel=9
\pdfadjustspacing=1 
\PassOptionsToPackage{pdftex}{graphicx}
    \usepackage{graphicx} 
\usepackage[type=CC,modifier=by-nc-nd,version=4.0,imagewidth=25mm]{doclicense}
\renewcommand*{\doclicenseURL}{https://creativecommons.org/licenses/by-nc-nd/4.0/}
\newcommand*{\myLicenseNotice}{Except where otherwise indicated, the original
    text and illustrations in \textit{\myTitle} by \myName{} and contributors
    are licensed under the \doclicenseLongName{} License}
\renewcommand*{\doclicenseLongTextForHyperref}{\myLicenseNotice.}

\newenvironment{spacedparagraphs}{%
    \par
    \setlength{\parindent}{0pt}%
    \setlength{\parskip}{0.75\baselineskip}%
}{\par}

% ********************************************************************
% Hyperreferences
% ********************************************************************
\hypersetup{%
    %draft, % = no hyperlinking at all (useful in b/w printouts)
    colorlinks=true, linktocpage=true, pdfstartpage=3, pdfstartview={FitV null},%
    % uncomment the following line if you want to have black links (e.g., for printing)
    %colorlinks=false, linktocpage=false, pdfstartpage=3, pdfstartview=FitV, pdfborder={0 0 0},%
    breaklinks=true, pdfpagemode=UseNone, pageanchor=true, pdfpagemode=UseOutlines,%
    plainpages=false, bookmarksnumbered, bookmarksopen=true, bookmarksopenlevel=1,%
    hypertexnames=true, pdfhighlight=/O,%nesting=true,%frenchlinks,%
    urlcolor=webbrown, linkcolor=RoyalBlue, citecolor=webgreen, %pagecolor=RoyalBlue,%
    %urlcolor=Black, linkcolor=Black, citecolor=Black, %pagecolor=Black,%
    keeppdfinfo, % Include the author in both PDF properties and XMP metadata.
    pdftitle={\myTitle},%
    pdfauthor={\myName},%
    pdfsubject={},%
    pdfkeywords={},%
    pdfcreator={pdfLaTeX},%
    pdfproducer={LaTeX with hyperref and classicthesis}%
}   

% ********************************************************************
% Setup autoreferences
% ********************************************************************
% There are some issues regarding autorefnames
% http://www.ureader.de/msg/136221647.aspx
% http://www.tex.ac.uk/cgi-bin/texfaq2html?label=latexwords
% you have to redefine the makros for the 
% language you use, e.g., american, ngerman
% (as chosen when loading babel/AtBeginDocument)
% ********************************************************************
\makeatletter
\@ifpackageloaded{babel}%
    {%
       \addto\extrasamerican{%
			\renewcommand*{\figureautorefname}{Figure}%
			\renewcommand*{\tableautorefname}{Table}%
			\renewcommand*{\partautorefname}{Part}%
			\renewcommand*{\chapterautorefname}{Chapter}%
			\renewcommand*{\sectionautorefname}{Section}%
			\renewcommand*{\subsectionautorefname}{Section}%
			\renewcommand*{\subsubsectionautorefname}{Section}%     
                }%
       \addto\extrasngerman{% 
			\renewcommand*{\paragraphautorefname}{Absatz}%
			\renewcommand*{\subparagraphautorefname}{Unterabsatz}%
			\renewcommand*{\footnoteautorefname}{Fu\"snote}%
			\renewcommand*{\FancyVerbLineautorefname}{Zeile}%
			\renewcommand*{\theoremautorefname}{Theorem}%
			\renewcommand*{\appendixautorefname}{Anhang}%
			\renewcommand*{\equationautorefname}{Gleichung}%        
			\renewcommand*{\itemautorefname}{Punkt}%
                }%  
            % Fix to getting autorefs for subfigures right (thanks to Belinda Vogt for changing the definition)
            \providecommand{\subfigureautorefname}{\figureautorefname}%             
    }{\relax}
\makeatother


% ****************************************************************************************************
% 7. Last calls before the bar closes
% ****************************************************************************************************

% ********************************************************************
% Development Stuff
% ********************************************************************
\listfiles
%\PassOptionsToPackage{l2tabu,orthodox,abort}{nag}
%   \usepackage{nag}
%\PassOptionsToPackage{warning, all}{onlyamsmath}
%   \usepackage{onlyamsmath}

% ********************************************************************
% Last, but not least...
% ********************************************************************
\usepackage{microtype} % Improves typography, reduces overflow via character protrusion/expansion
\usepackage[dottedtoc]{classicthesis} 
\usepackage{lmodern}

% ****************************************************************************************************

% ****************************************************************************************************
% 8. Further adjustments (experimental)
% ****************************************************************************************************
% ********************************************************************
% Changing the text area
% ********************************************************************
\linespread{1.05} % a bit more for Palatino
% \areaset[current]{312pt}{761pt} % 686 (factor 2.2) + 33 head + 42 head \the\footskip
% \setlength{\marginparwidth}{7em}%
% \setlength{\marginparsep}{2em}%

% ********************************************************************
% Fix text/equation overflow into margins
% ********************************************************************
\tolerance=1500                    % Allow looser spacing (default 200)
\emergencystretch=1.5em            % Extra stretch to avoid overflow
\hfuzz=1.5pt                       % Suppress warnings for minor overflow
\hbadness=1500   

% ********************************************************************
% Using different fonts
% ********************************************************************
%\usepackage[oldstylenums]{kpfonts} % oldstyle notextcomp
%\usepackage[osf]{libertine}
%\usepackage[light,condensed,math]{iwona}
%\renewcommand{\sfdefault}{iwona}
%\usepackage{lmodern} % <-- no osf support :-(
%\usepackage{cfr-lm} % 
%\usepackage[urw-garamond]{mathdesign} <-- no osf support :-(
%\usepackage[default,osfigures]{opensans} % scale=0.95 
%\usepackage[sfdefault]{FiraSans}
% ****************************************************************************************************

% ****************************************************************************************************
% X. My Custom Stuff
% ****************************************************************************************************
\usepackage{nomencl}
\usepackage{setspace}        % for \onehalfspacing
\usepackage{svg}
\usepackage{epigraph}
\usepackage{float}
\usepackage[section]{placeins} % keep floats within their sections/chapters
\usepackage{booktabs}
\usepackage{import}
\usepackage{array}
\usepackage{ragged2e}
\usepackage{longtable}
\usepackage{multirow}
\setlength{\LTleft}{0pt} % Fix longtable shift on page breaks in ClassicThesis
\setlength{\LTright}{\fill}
\usepackage{tikz}
\usepackage{lipsum}
% \usepackage[scaled=0.85]{beramono} % reload, after lmodern
% \renewcommand{\ttdefault}{fvm}     % make sure \ttfamily = Bera Mono
\usepackage[table]{xcolor}
\usepackage[most]{tcolorbox}
\usepackage[para]{threeparttable}
\usepackage{adjustbox}

\usepackage{xkeyval}     % for \define@key[FigTools]{...}, \setkeys
\usepackage{ifthen}      % for \ifthenelse and \equal
\usepackage{ifoddpage}   % for \checkoddpage and \ifoddpage

% Algorithms
\usepackage{algorithm}
\usepackage{algpseudocode}
% Each algorithm restarts its line numbers; include the float in PDF destinations.
\counterwithin*{ALG@line}{algorithm}

\usepackage{url}

% ********************************************************************
% CLEVEREF SETUP
% ********************************************************************
\usepackage[nameinlink,capitalize]{cleveref}
% Older LaTeX formats lose all but the first letter of acronym labels through
% cleveref's amsmath adapter. Use the complete key for acronym first-use links.
% https://github.com/oetiker/acronym/issues/62
\makeatletter
\IfFormatAtLeastTF{2025-06-01}{}{%
    \patchcmd{\AC@verridelabel}
        {\ifdefined\ltx@label\ltx@label}
        {\ifdefined\cref@label\cref@label}{}{}%
}
\makeatother

% 1. GENERIC SETTINGS
% -------------------
% Capitalize all references (even mid-sentence) to match Engineering standards
% e.g., "as seen in Fig. 1" rather than "as seen in fig. 1"

% 2. SPECIFIC DEFINITIONS
% -----------------------

% Equations: "Eq." and "Eqs."
\crefname{equation}{Eq.}{Eqs.}
\Crefname{equation}{Eq.}{Eqs.}
% Optional: Remove parentheses around eq number (e.g., "Eq. 1" instead of "Eq. (1)")
% If you prefer standard math format "Eq. (1)", comment out the line below.
\creflabelformat{equation}{#2#1#3}

% Figures: "Fig." and "Figs."
\crefname{figure}{Fig.}{Figs.}
\Crefname{figure}{Fig.}{Figs.}

% Tables: "Table" and "Tables" (Usually fully written out)
\crefname{table}{Table}{Tables}
\Crefname{table}{Table}{Tables}

% Code Listings (from listings package): "Listing" and "Listings"
\crefname{lstlisting}{Listing}{Listings}
\Crefname{lstlisting}{Listing}{Listings}

% Algorithms (from algorithm package): "Algorithm" and "Algorithms"
\crefname{algorithm}{Algorithm}{Algorithms}
\Crefname{algorithm}{Algorithm}{Algorithms}

% Sections/Chapters: "Section", "Chapter", "Part"
\crefname{part}{Part}{Parts}
\Crefname{part}{Part}{Parts}
\renewcommand{\thepart}{\Roman{part}}
% A ClassicThesis part ships its title page before \part returns.
\makeatletter
\newcommand{\labeledpart}[2]{%
    \GetTitleString{#1}\let\@currentlabelname\GetTitleStringResult
    \part[#1]{#1\label{#2}}%
}
\makeatother

\crefname{chapter}{Chapter}{Chapters}
\Crefname{chapter}{Chapter}{Chapters}

\crefname{section}{Section}{Sections}
\Crefname{section}{Section}{Sections}

\crefname{subsection}{Section}{Sections} % Map subsections to "Section"
\Crefname{subsection}{Section}{Sections}

\crefname{subsubsection}{Section}{Sections} % Map sub-subsections to "Section"
\Crefname{subsubsection}{Section}{Sections}

% Appendices
\crefname{appendix}{Appendix}{Appendices}
\Crefname{appendix}{Appendix}{Appendices}

% Enumerated Lists (items)
\crefname{enumi}{Item}{Items}
\Crefname{enumi}{Item}{Items}
\crefname{enumii}{Item}{Items}
\Crefname{enumii}{Item}{Items}
\crefname{enumiii}{Item}{Items}
\Crefname{enumiii}{Item}{Items}

% url penalties
\setcounter{biburllcpenalty}{7000}
\setcounter{biburlucpenalty}{8000}

% your column types
\newcolumntype{L}[1]{>{\RaggedRight\arraybackslash}p{#1}}
\newcolumntype{R}[1]{>{\RaggedLeft\arraybackslash}p{#1}}

% tcolorbox style
% Define color themes
\definecolor{lucasredlight}{RGB}{255,230,230}
\definecolor{lucasreddark}{RGB}{180,0,0}

\definecolor{lucasbluelight}{RGB}{225,235,255}
\definecolor{lucasbluedark}{RGB}{0,60,150}

% A key to choose theme = red or blue
\tcbset{
  lucasbox/.style={
    enhanced,
    fonttitle=\bfseries,
    title filled,
    boxrule=1pt,
    toptitle=2mm,
    bottomtitle=1mm
  },
  lucasboxtheme/red/.style={
    colback=lucasredlight,
    colframe=lucasreddark,
    coltitle=white,
    colbacktitle=lucasreddark
  },
  lucasboxtheme/blue/.style={
    colback=lucasbluelight,
    colframe=lucasbluedark,
    coltitle=white,
    colbacktitle=lucasbluedark
  }
}

% Default theme = red
\newenvironment{importantbox}[2][red]
  {\begin{tcolorbox}[lucasbox, lucasboxtheme/#1, title=#2]}
  {\end{tcolorbox}}

\DeclareMathOperator*{\argmax}{arg\,max}

\newcommand{\quotegraffito}[2]{\graffito{#1\\--~\emph{#2}}}

% grey midrule for trial block boundaries
\newcommand{\graymidrule}{%
  \arrayrulecolor{gray!60}%
  \midrule
  \arrayrulecolor{black}%
}

% Method labels (math and text safe)
\newcommand{\mBase}{\text{\textsc{Base}}}   % pretrained / SFT base policy
\newcommand{\mRL}{\text{\textsc{RL}}}       % pure RL from scratch
\newcommand{\mResRL}{\text{\textsc{ResRL}}} % residual RL on Base

% ****************************************************************************************************  
% If you like the classicthesis, then I would appreciate a postcard. 
% My address can be found in the file ClassicThesis.pdf. A collection 
% of the postcards I received so far is available online at 
% http://postcards.miede.de
% ****************************************************************************************************


% ****************************************************************************************************
% 0. Set the encoding of your files. UTF-8 is the only sensible encoding nowadays. If you can't read
% äöüßáéçèê∂åëæƒÏ€ then change the encoding setting in your editor, not the line below. If your editor
% does not support utf8 use another editor!
% ****************************************************************************************************
\PassOptionsToPackage{utf8}{inputenc}
	\usepackage{inputenc}

% ****************************************************************************************************
% 1. Configure classicthesis for your needs here, e.g., remove "drafting" below 
% in order to deactivate the time-stamp on the pages
% ****************************************************************************************************
\PassOptionsToPackage{eulerchapternumbers,listings,%
					 pdfspacing,%floatperchapter,%linedheaders,%
					 subfig,beramono,parts}{classicthesis}                                        
% ********************************************************************
% Available options for classicthesis.sty 
% (see ClassicThesis.pdf for more information):
% drafting
% parts nochapters linedheaders
% eulerchapternumbers beramono eulermath pdfspacing minionprospacing
% tocaligned dottedtoc manychapters
% listings floatperchapter subfig
% ********************************************************************


% ****************************************************************************************************
% 2. Personal data and user ad-hoc commands
% ****************************************************************************************************
\newcommand{\myTitle}{Real-World Reinforcement Learning for Vision-Language Robotic Manipulation\xspace}
\newcommand{\mySubtitle}{Deep-learning for Robotics\xspace}
\newcommand{\myDegree}{Bachelor of Engineering Honours and Computer Science\xspace}
\newcommand{\myName}{Lucas Barbosa\xspace}
\newcommand{\myProf}{Will Midgley\xspace}
\newcommand{\myOtherProf}{Put name here\xspace}
\newcommand{\mySupervisor}{Will Midgley\xspace}
\newcommand{\myFaculty}{Faculty of Engineering\xspace}
\newcommand{\myDepartment}{School of Mechanical and Manufacturing Engineering\xspace}
\newcommand{\myUni}{UNSW\xspace}
\newcommand{\myLocation}{Sydney\xspace}
% \newcommand{\myTime}{November 2025\xspace}
% \newcommand{\myVersion}{version 4.2\xspace}

% ********************************************************************
% Setup, finetuning, and useful commands
% ********************************************************************
\newlength{\abcd} % for ab..z string length calculation
\providecommand{\mLyX}{L\kern-.1667em\lower.25em\hbox{Y}\kern-.125emX\@}
\newcommand{\ie}{i.\,e.}
\newcommand{\Ie}{I.\,e.}
\newcommand{\eg}{e.\,g.}
\newcommand{\Eg}{E.\,g.} 
% ****************************************************************************************************


% ****************************************************************************************************
% 3. Loading some handy packages
% ****************************************************************************************************
% ******************************************************************** 
% Packages with options that might require adjustments
% ******************************************************************** 
%\PassOptionsToPackage{ngerman,american}{babel}   % change this to your language(s)
% Spanish languages need extra options in order to work with this template
%\PassOptionsToPackage{spanish,es-lcroman}{babel}
	\usepackage{babel}                  

\usepackage{csquotes}
\PassOptionsToPackage{%
    backend=biber,%
	language=auto,%
	style=numeric-comp,%
    %style=authoryear-comp, % Author 1999, 2010
    %bibstyle=authoryear,dashed=false, % dashed: substitute rep. author with ---
    sorting=nyt, % name, year, title
    maxbibnames=10, % default: 3, et al.
    %backref=true,%
    natbib=true % natbib compatibility mode (\citep and \citet still work)
}{biblatex}
    \usepackage{biblatex}

% Maths
\usepackage{amsmath, amssymb}

% ******************************************************************** 
% General useful packages
% ******************************************************************** 
\PassOptionsToPackage{T1}{fontenc} % T2A for cyrillics
    \usepackage{fontenc}     
\usepackage{textcomp} % fix warning with missing font shapes
\usepackage{scrhack} % fix warnings when using KOMA with listings package          
\usepackage{xspace} % to get the spacing after macros right  
\usepackage{mparhack} % get marginpar right
%\usepackage[latest]{latexrelease} % will be used once available in more distributions (ISSUE #107)
\PassOptionsToPackage{printonlyused,smaller}{acronym} 
    \usepackage{acronym} % nice macros for handling all acronyms in the thesis
    %\renewcommand{\bflabel}[1]{{#1}\hfill} % fix the list of acronyms --> no longer working
    %\renewcommand*{\acsfont}[1]{\textsc{#1}} 
    \renewcommand*{\aclabelfont}[1]{\acsfont{#1}}
% ****************************************************************************************************


% ****************************************************************************************************
% 4. Setup floats: tables, (sub)figures, and captions
% ****************************************************************************************************
\usepackage{tabularx} % better tables
    \setlength{\extrarowheight}{3pt} % increase table row height
\newcommand{\tableheadline}[1]{\multicolumn{1}{c}{\spacedlowsmallcaps{#1}}}
\newcommand{\myfloatalign}{\centering} % to be used with each float for alignment
% \usepackage{caption}
% \usepackage[labelsep=period, labelformat=default,font=footnotesize,labelfont=bf, width=\textwidth,format=plain]{caption}
% Thanks to cgnieder and Claus Lahiri
% http://tex.stackexchange.com/questions/69349/spacedlowsmallcaps-in-caption-label
% [REMOVED DUE TO OTHER PROBLEMS, SEE ISSUE #82]    
%\DeclareCaptionLabelFormat{smallcaps}{\bothIfFirst{#1}{~}\MakeTextLowercase{\textsc{#2}}}
%\captionsetup{font=small,labelformat=smallcaps} % format=hang,
\usepackage[labelsep=period,labelformat=default,font=footnotesize,labelfont=bf,format=plain]{caption}
\usepackage{subfig}
% ****************************************************************************************************


% ****************************************************************************************************
% 5. Setup code listings
% ****************************************************************************************************
\usepackage{listings} 
%\lstset{emph={trueIndex,root},emphstyle=\color{BlueViolet}}%\underbar} % for special keywords
\lstset{language=[LaTeX]Tex,%C++,
    morekeywords={PassOptionsToPackage,selectlanguage},
    keywordstyle=\color{RoyalBlue},%\bfseries,
    basicstyle=\small\ttfamily,
    %identifierstyle=\color{NavyBlue},
    commentstyle=\color{Green}\ttfamily,
    stringstyle=\rmfamily,
    numbers=none,%left,%
    numberstyle=\scriptsize,%\tiny
    stepnumber=5,
    numbersep=8pt,
    showstringspaces=false,
    breaklines=true,
    %frameround=ftff,
    %frame=single,
    belowcaptionskip=.75\baselineskip
    %frame=L
} 
% ****************************************************************************************************             

% ****************************************************************************************************
% 6. PDFLaTeX, hyperreferences and citation backreferences
% ****************************************************************************************************
% ********************************************************************
% Using PDFLaTeX
% ********************************************************************
\usepackage{titlesec} % Let hyperref place anchors with the typeset headings.
\PassOptionsToPackage{pdftex,hyperfootnotes=false,pdfpagelabels}{hyperref}
    \usepackage{hyperref}  % backref linktocpage pagebackref
\pdfcompresslevel=9
\pdfadjustspacing=1 
\PassOptionsToPackage{pdftex}{graphicx}
    \usepackage{graphicx} 
\usepackage[type=CC,modifier=by-nc-nd,version=4.0,imagewidth=25mm]{doclicense}
\renewcommand*{\doclicenseURL}{https://creativecommons.org/licenses/by-nc-nd/4.0/}
\newcommand*{\myLicenseNotice}{Except where otherwise indicated, the original
    text and illustrations in \textit{\myTitle} by \myName{} and contributors
    are licensed under the \doclicenseLongName{} License}
\renewcommand*{\doclicenseLongTextForHyperref}{\myLicenseNotice.}

\newenvironment{spacedparagraphs}{%
    \par
    \setlength{\parindent}{0pt}%
    \setlength{\parskip}{0.75\baselineskip}%
}{\par}

% ********************************************************************
% Hyperreferences
% ********************************************************************
\hypersetup{%
    %draft, % = no hyperlinking at all (useful in b/w printouts)
    colorlinks=true, linktocpage=true, pdfstartpage=3, pdfstartview={FitV null},%
    % uncomment the following line if you want to have black links (e.g., for printing)
    %colorlinks=false, linktocpage=false, pdfstartpage=3, pdfstartview=FitV, pdfborder={0 0 0},%
    breaklinks=true, pdfpagemode=UseNone, pageanchor=true, pdfpagemode=UseOutlines,%
    plainpages=false, bookmarksnumbered, bookmarksopen=true, bookmarksopenlevel=1,%
    hypertexnames=true, pdfhighlight=/O,%nesting=true,%frenchlinks,%
    urlcolor=webbrown, linkcolor=RoyalBlue, citecolor=webgreen, %pagecolor=RoyalBlue,%
    %urlcolor=Black, linkcolor=Black, citecolor=Black, %pagecolor=Black,%
    keeppdfinfo, % Include the author in both PDF properties and XMP metadata.
    pdftitle={\myTitle},%
    pdfauthor={\myName},%
    pdfsubject={},%
    pdfkeywords={},%
    pdfcreator={pdfLaTeX},%
    pdfproducer={LaTeX with hyperref and classicthesis}%
}   

% ********************************************************************
% Setup autoreferences
% ********************************************************************
% There are some issues regarding autorefnames
% http://www.ureader.de/msg/136221647.aspx
% http://www.tex.ac.uk/cgi-bin/texfaq2html?label=latexwords
% you have to redefine the makros for the 
% language you use, e.g., american, ngerman
% (as chosen when loading babel/AtBeginDocument)
% ********************************************************************
\makeatletter
\@ifpackageloaded{babel}%
    {%
       \addto\extrasamerican{%
			\renewcommand*{\figureautorefname}{Figure}%
			\renewcommand*{\tableautorefname}{Table}%
			\renewcommand*{\partautorefname}{Part}%
			\renewcommand*{\chapterautorefname}{Chapter}%
			\renewcommand*{\sectionautorefname}{Section}%
			\renewcommand*{\subsectionautorefname}{Section}%
			\renewcommand*{\subsubsectionautorefname}{Section}%     
                }%
       \addto\extrasngerman{% 
			\renewcommand*{\paragraphautorefname}{Absatz}%
			\renewcommand*{\subparagraphautorefname}{Unterabsatz}%
			\renewcommand*{\footnoteautorefname}{Fu\"snote}%
			\renewcommand*{\FancyVerbLineautorefname}{Zeile}%
			\renewcommand*{\theoremautorefname}{Theorem}%
			\renewcommand*{\appendixautorefname}{Anhang}%
			\renewcommand*{\equationautorefname}{Gleichung}%        
			\renewcommand*{\itemautorefname}{Punkt}%
                }%  
            % Fix to getting autorefs for subfigures right (thanks to Belinda Vogt for changing the definition)
            \providecommand{\subfigureautorefname}{\figureautorefname}%             
    }{\relax}
\makeatother


% ****************************************************************************************************
% 7. Last calls before the bar closes
% ****************************************************************************************************

% ********************************************************************
% Development Stuff
% ********************************************************************
\listfiles
%\PassOptionsToPackage{l2tabu,orthodox,abort}{nag}
%   \usepackage{nag}
%\PassOptionsToPackage{warning, all}{onlyamsmath}
%   \usepackage{onlyamsmath}

% ********************************************************************
% Last, but not least...
% ********************************************************************
\usepackage{microtype} % Improves typography, reduces overflow via character protrusion/expansion
\usepackage[dottedtoc]{classicthesis} 
\usepackage{lmodern}

% ****************************************************************************************************

% ****************************************************************************************************
% 8. Further adjustments (experimental)
% ****************************************************************************************************
% ********************************************************************
% Changing the text area
% ********************************************************************
\linespread{1.05} % a bit more for Palatino
% \areaset[current]{312pt}{761pt} % 686 (factor 2.2) + 33 head + 42 head \the\footskip
% \setlength{\marginparwidth}{7em}%
% \setlength{\marginparsep}{2em}%

% ********************************************************************
% Fix text/equation overflow into margins
% ********************************************************************
\tolerance=1500                    % Allow looser spacing (default 200)
\emergencystretch=1.5em            % Extra stretch to avoid overflow
\hfuzz=1.5pt                       % Suppress warnings for minor overflow
\hbadness=1500   

% ********************************************************************
% Using different fonts
% ********************************************************************
%\usepackage[oldstylenums]{kpfonts} % oldstyle notextcomp
%\usepackage[osf]{libertine}
%\usepackage[light,condensed,math]{iwona}
%\renewcommand{\sfdefault}{iwona}
%\usepackage{lmodern} % <-- no osf support :-(
%\usepackage{cfr-lm} % 
%\usepackage[urw-garamond]{mathdesign} <-- no osf support :-(
%\usepackage[default,osfigures]{opensans} % scale=0.95 
%\usepackage[sfdefault]{FiraSans}
% ****************************************************************************************************

% ****************************************************************************************************
% X. My Custom Stuff
% ****************************************************************************************************
\usepackage{nomencl}
\usepackage{setspace}        % for \onehalfspacing
\usepackage{svg}
\usepackage{epigraph}
\usepackage{float}
\usepackage[section]{placeins} % keep floats within their sections/chapters
\usepackage{booktabs}
\usepackage{import}
\usepackage{array}
\usepackage{ragged2e}
\usepackage{longtable}
\usepackage{multirow}
\setlength{\LTleft}{0pt} % Fix longtable shift on page breaks in ClassicThesis
\setlength{\LTright}{\fill}
\usepackage{tikz}
\usepackage{lipsum}
% \usepackage[scaled=0.85]{beramono} % reload, after lmodern
% \renewcommand{\ttdefault}{fvm}     % make sure \ttfamily = Bera Mono
\usepackage[table]{xcolor}
\usepackage[most]{tcolorbox}
\usepackage[para]{threeparttable}
\usepackage{adjustbox}

\usepackage{xkeyval}     % for \define@key[FigTools]{...}, \setkeys
\usepackage{ifthen}      % for \ifthenelse and \equal
\usepackage{ifoddpage}   % for \checkoddpage and \ifoddpage

% Algorithms
\usepackage{algorithm}
\usepackage{algpseudocode}
% Each algorithm restarts its line numbers; include the float in PDF destinations.
\counterwithin*{ALG@line}{algorithm}

\usepackage{url}

% ********************************************************************
% CLEVEREF SETUP
% ********************************************************************
\usepackage[nameinlink,capitalize]{cleveref}
% Older LaTeX formats lose all but the first letter of acronym labels through
% cleveref's amsmath adapter. Use the complete key for acronym first-use links.
% https://github.com/oetiker/acronym/issues/62
\makeatletter
\IfFormatAtLeastTF{2025-06-01}{}{%
    \patchcmd{\AC@verridelabel}
        {\ifdefined\ltx@label\ltx@label}
        {\ifdefined\cref@label\cref@label}{}{}%
}
\makeatother

% 1. GENERIC SETTINGS
% -------------------
% Capitalize all references (even mid-sentence) to match Engineering standards
% e.g., "as seen in Fig. 1" rather than "as seen in fig. 1"

% 2. SPECIFIC DEFINITIONS
% -----------------------

% Equations: "Eq." and "Eqs."
\crefname{equation}{Eq.}{Eqs.}
\Crefname{equation}{Eq.}{Eqs.}
% Optional: Remove parentheses around eq number (e.g., "Eq. 1" instead of "Eq. (1)")
% If you prefer standard math format "Eq. (1)", comment out the line below.
\creflabelformat{equation}{#2#1#3}

% Figures: "Fig." and "Figs."
\crefname{figure}{Fig.}{Figs.}
\Crefname{figure}{Fig.}{Figs.}

% Tables: "Table" and "Tables" (Usually fully written out)
\crefname{table}{Table}{Tables}
\Crefname{table}{Table}{Tables}

% Code Listings (from listings package): "Listing" and "Listings"
\crefname{lstlisting}{Listing}{Listings}
\Crefname{lstlisting}{Listing}{Listings}

% Algorithms (from algorithm package): "Algorithm" and "Algorithms"
\crefname{algorithm}{Algorithm}{Algorithms}
\Crefname{algorithm}{Algorithm}{Algorithms}

% Sections/Chapters: "Section", "Chapter", "Part"
\crefname{part}{Part}{Parts}
\Crefname{part}{Part}{Parts}
\renewcommand{\thepart}{\Roman{part}}
% A ClassicThesis part ships its title page before \part returns.
\makeatletter
\newcommand{\labeledpart}[2]{%
    \GetTitleString{#1}\let\@currentlabelname\GetTitleStringResult
    \part[#1]{#1\label{#2}}%
}
\makeatother

\crefname{chapter}{Chapter}{Chapters}
\Crefname{chapter}{Chapter}{Chapters}

\crefname{section}{Section}{Sections}
\Crefname{section}{Section}{Sections}

\crefname{subsection}{Section}{Sections} % Map subsections to "Section"
\Crefname{subsection}{Section}{Sections}

\crefname{subsubsection}{Section}{Sections} % Map sub-subsections to "Section"
\Crefname{subsubsection}{Section}{Sections}

% Appendices
\crefname{appendix}{Appendix}{Appendices}
\Crefname{appendix}{Appendix}{Appendices}

% Enumerated Lists (items)
\crefname{enumi}{Item}{Items}
\Crefname{enumi}{Item}{Items}
\crefname{enumii}{Item}{Items}
\Crefname{enumii}{Item}{Items}
\crefname{enumiii}{Item}{Items}
\Crefname{enumiii}{Item}{Items}

% url penalties
\setcounter{biburllcpenalty}{7000}
\setcounter{biburlucpenalty}{8000}

% your column types
\newcolumntype{L}[1]{>{\RaggedRight\arraybackslash}p{#1}}
\newcolumntype{R}[1]{>{\RaggedLeft\arraybackslash}p{#1}}

% tcolorbox style
% Define color themes
\definecolor{lucasredlight}{RGB}{255,230,230}
\definecolor{lucasreddark}{RGB}{180,0,0}

\definecolor{lucasbluelight}{RGB}{225,235,255}
\definecolor{lucasbluedark}{RGB}{0,60,150}

% A key to choose theme = red or blue
\tcbset{
  lucasbox/.style={
    enhanced,
    fonttitle=\bfseries,
    title filled,
    boxrule=1pt,
    toptitle=2mm,
    bottomtitle=1mm
  },
  lucasboxtheme/red/.style={
    colback=lucasredlight,
    colframe=lucasreddark,
    coltitle=white,
    colbacktitle=lucasreddark
  },
  lucasboxtheme/blue/.style={
    colback=lucasbluelight,
    colframe=lucasbluedark,
    coltitle=white,
    colbacktitle=lucasbluedark
  }
}

% Default theme = red
\newenvironment{importantbox}[2][red]
  {\begin{tcolorbox}[lucasbox, lucasboxtheme/#1, title=#2]}
  {\end{tcolorbox}}

\DeclareMathOperator*{\argmax}{arg\,max}

\newcommand{\quotegraffito}[2]{\graffito{#1\\--~\emph{#2}}}

% grey midrule for trial block boundaries
\newcommand{\graymidrule}{%
  \arrayrulecolor{gray!60}%
  \midrule
  \arrayrulecolor{black}%
}

% Method labels (math and text safe)
\newcommand{\mBase}{\text{\textsc{Base}}}   % pretrained / SFT base policy
\newcommand{\mRL}{\text{\textsc{RL}}}       % pure RL from scratch
\newcommand{\mResRL}{\text{\textsc{ResRL}}} % residual RL on Base

% ****************************************************************************************************  
% If you like the classicthesis, then I would appreciate a postcard. 
% My address can be found in the file ClassicThesis.pdf. A collection 
% of the postcards I received so far is available online at 
% http://postcards.miede.de
% ****************************************************************************************************


% ****************************************************************************************************
% 0. Set the encoding of your files. UTF-8 is the only sensible encoding nowadays. If you can't read
% äöüßáéçèê∂åëæƒÏ€ then change the encoding setting in your editor, not the line below. If your editor
% does not support utf8 use another editor!
% ****************************************************************************************************
\PassOptionsToPackage{utf8}{inputenc}
	\usepackage{inputenc}

% ****************************************************************************************************
% 1. Configure classicthesis for your needs here, e.g., remove "drafting" below 
% in order to deactivate the time-stamp on the pages
% ****************************************************************************************************
\PassOptionsToPackage{eulerchapternumbers,listings,%
					 pdfspacing,%floatperchapter,%linedheaders,%
					 subfig,beramono,parts}{classicthesis}                                        
% ********************************************************************
% Available options for classicthesis.sty 
% (see ClassicThesis.pdf for more information):
% drafting
% parts nochapters linedheaders
% eulerchapternumbers beramono eulermath pdfspacing minionprospacing
% tocaligned dottedtoc manychapters
% listings floatperchapter subfig
% ********************************************************************


% ****************************************************************************************************
% 2. Personal data and user ad-hoc commands
% ****************************************************************************************************
\newcommand{\myTitle}{Real-World Reinforcement Learning for Vision-Language Robotic Manipulation\xspace}
\newcommand{\mySubtitle}{Deep-learning for Robotics\xspace}
\newcommand{\myDegree}{Bachelor of Engineering Honours and Computer Science\xspace}
\newcommand{\myName}{Lucas Barbosa\xspace}
\newcommand{\myProf}{Will Midgley\xspace}
\newcommand{\myOtherProf}{Put name here\xspace}
\newcommand{\mySupervisor}{Will Midgley\xspace}
\newcommand{\myFaculty}{Faculty of Engineering\xspace}
\newcommand{\myDepartment}{School of Mechanical and Manufacturing Engineering\xspace}
\newcommand{\myUni}{UNSW\xspace}
\newcommand{\myLocation}{Sydney\xspace}
% \newcommand{\myTime}{November 2025\xspace}
% \newcommand{\myVersion}{version 4.2\xspace}

% ********************************************************************
% Setup, finetuning, and useful commands
% ********************************************************************
\newlength{\abcd} % for ab..z string length calculation
\providecommand{\mLyX}{L\kern-.1667em\lower.25em\hbox{Y}\kern-.125emX\@}
\newcommand{\ie}{i.\,e.}
\newcommand{\Ie}{I.\,e.}
\newcommand{\eg}{e.\,g.}
\newcommand{\Eg}{E.\,g.} 
% ****************************************************************************************************


% ****************************************************************************************************
% 3. Loading some handy packages
% ****************************************************************************************************
% ******************************************************************** 
% Packages with options that might require adjustments
% ******************************************************************** 
%\PassOptionsToPackage{ngerman,american}{babel}   % change this to your language(s)
% Spanish languages need extra options in order to work with this template
%\PassOptionsToPackage{spanish,es-lcroman}{babel}
	\usepackage{babel}                  

\usepackage{csquotes}
\PassOptionsToPackage{%
    backend=biber,%
	language=auto,%
	style=numeric-comp,%
    %style=authoryear-comp, % Author 1999, 2010
    %bibstyle=authoryear,dashed=false, % dashed: substitute rep. author with ---
    sorting=nyt, % name, year, title
    maxbibnames=10, % default: 3, et al.
    %backref=true,%
    natbib=true % natbib compatibility mode (\citep and \citet still work)
}{biblatex}
    \usepackage{biblatex}

% Maths
\usepackage{amsmath, amssymb}

% ******************************************************************** 
% General useful packages
% ******************************************************************** 
\PassOptionsToPackage{T1}{fontenc} % T2A for cyrillics
    \usepackage{fontenc}     
\usepackage{textcomp} % fix warning with missing font shapes
\usepackage{scrhack} % fix warnings when using KOMA with listings package          
\usepackage{xspace} % to get the spacing after macros right  
\usepackage{mparhack} % get marginpar right
%\usepackage[latest]{latexrelease} % will be used once available in more distributions (ISSUE #107)
\PassOptionsToPackage{printonlyused,smaller}{acronym} 
    \usepackage{acronym} % nice macros for handling all acronyms in the thesis
    %\renewcommand{\bflabel}[1]{{#1}\hfill} % fix the list of acronyms --> no longer working
    %\renewcommand*{\acsfont}[1]{\textsc{#1}} 
    \renewcommand*{\aclabelfont}[1]{\acsfont{#1}}
% ****************************************************************************************************


% ****************************************************************************************************
% 4. Setup floats: tables, (sub)figures, and captions
% ****************************************************************************************************
\usepackage{tabularx} % better tables
    \setlength{\extrarowheight}{3pt} % increase table row height
\newcommand{\tableheadline}[1]{\multicolumn{1}{c}{\spacedlowsmallcaps{#1}}}
\newcommand{\myfloatalign}{\centering} % to be used with each float for alignment
% \usepackage{caption}
% \usepackage[labelsep=period, labelformat=default,font=footnotesize,labelfont=bf, width=\textwidth,format=plain]{caption}
% Thanks to cgnieder and Claus Lahiri
% http://tex.stackexchange.com/questions/69349/spacedlowsmallcaps-in-caption-label
% [REMOVED DUE TO OTHER PROBLEMS, SEE ISSUE #82]    
%\DeclareCaptionLabelFormat{smallcaps}{\bothIfFirst{#1}{~}\MakeTextLowercase{\textsc{#2}}}
%\captionsetup{font=small,labelformat=smallcaps} % format=hang,
\usepackage[labelsep=period,labelformat=default,font=footnotesize,labelfont=bf,format=plain]{caption}
\usepackage{subfig}
% ****************************************************************************************************


% ****************************************************************************************************
% 5. Setup code listings
% ****************************************************************************************************
\usepackage{listings} 
%\lstset{emph={trueIndex,root},emphstyle=\color{BlueViolet}}%\underbar} % for special keywords
\lstset{language=[LaTeX]Tex,%C++,
    morekeywords={PassOptionsToPackage,selectlanguage},
    keywordstyle=\color{RoyalBlue},%\bfseries,
    basicstyle=\small\ttfamily,
    %identifierstyle=\color{NavyBlue},
    commentstyle=\color{Green}\ttfamily,
    stringstyle=\rmfamily,
    numbers=none,%left,%
    numberstyle=\scriptsize,%\tiny
    stepnumber=5,
    numbersep=8pt,
    showstringspaces=false,
    breaklines=true,
    %frameround=ftff,
    %frame=single,
    belowcaptionskip=.75\baselineskip
    %frame=L
} 
% ****************************************************************************************************             

% ****************************************************************************************************
% 6. PDFLaTeX, hyperreferences and citation backreferences
% ****************************************************************************************************
% ********************************************************************
% Using PDFLaTeX
% ********************************************************************
\usepackage{titlesec} % Let hyperref place anchors with the typeset headings.
\PassOptionsToPackage{pdftex,hyperfootnotes=false,pdfpagelabels}{hyperref}
    \usepackage{hyperref}  % backref linktocpage pagebackref
\pdfcompresslevel=9
\pdfadjustspacing=1 
\PassOptionsToPackage{pdftex}{graphicx}
    \usepackage{graphicx} 
\usepackage[type=CC,modifier=by-nc-nd,version=4.0,imagewidth=25mm]{doclicense}
\renewcommand*{\doclicenseURL}{https://creativecommons.org/licenses/by-nc-nd/4.0/}
\newcommand*{\myLicenseNotice}{Except where otherwise indicated, the original
    text and illustrations in \textit{\myTitle} by \myName{} and contributors
    are licensed under the \doclicenseLongName{} License}
\renewcommand*{\doclicenseLongTextForHyperref}{\myLicenseNotice.}

\newenvironment{spacedparagraphs}{%
    \par
    \setlength{\parindent}{0pt}%
    \setlength{\parskip}{0.75\baselineskip}%
}{\par}

% ********************************************************************
% Hyperreferences
% ********************************************************************
\hypersetup{%
    %draft, % = no hyperlinking at all (useful in b/w printouts)
    colorlinks=true, linktocpage=true, pdfstartpage=3, pdfstartview={FitV null},%
    % uncomment the following line if you want to have black links (e.g., for printing)
    %colorlinks=false, linktocpage=false, pdfstartpage=3, pdfstartview=FitV, pdfborder={0 0 0},%
    breaklinks=true, pdfpagemode=UseNone, pageanchor=true, pdfpagemode=UseOutlines,%
    plainpages=false, bookmarksnumbered, bookmarksopen=true, bookmarksopenlevel=1,%
    hypertexnames=true, pdfhighlight=/O,%nesting=true,%frenchlinks,%
    urlcolor=webbrown, linkcolor=RoyalBlue, citecolor=webgreen, %pagecolor=RoyalBlue,%
    %urlcolor=Black, linkcolor=Black, citecolor=Black, %pagecolor=Black,%
    keeppdfinfo, % Include the author in both PDF properties and XMP metadata.
    pdftitle={\myTitle},%
    pdfauthor={\myName},%
    pdfsubject={},%
    pdfkeywords={},%
    pdfcreator={pdfLaTeX},%
    pdfproducer={LaTeX with hyperref and classicthesis}%
}   

% ********************************************************************
% Setup autoreferences
% ********************************************************************
% There are some issues regarding autorefnames
% http://www.ureader.de/msg/136221647.aspx
% http://www.tex.ac.uk/cgi-bin/texfaq2html?label=latexwords
% you have to redefine the makros for the 
% language you use, e.g., american, ngerman
% (as chosen when loading babel/AtBeginDocument)
% ********************************************************************
\makeatletter
\@ifpackageloaded{babel}%
    {%
       \addto\extrasamerican{%
			\renewcommand*{\figureautorefname}{Figure}%
			\renewcommand*{\tableautorefname}{Table}%
			\renewcommand*{\partautorefname}{Part}%
			\renewcommand*{\chapterautorefname}{Chapter}%
			\renewcommand*{\sectionautorefname}{Section}%
			\renewcommand*{\subsectionautorefname}{Section}%
			\renewcommand*{\subsubsectionautorefname}{Section}%     
                }%
       \addto\extrasngerman{% 
			\renewcommand*{\paragraphautorefname}{Absatz}%
			\renewcommand*{\subparagraphautorefname}{Unterabsatz}%
			\renewcommand*{\footnoteautorefname}{Fu\"snote}%
			\renewcommand*{\FancyVerbLineautorefname}{Zeile}%
			\renewcommand*{\theoremautorefname}{Theorem}%
			\renewcommand*{\appendixautorefname}{Anhang}%
			\renewcommand*{\equationautorefname}{Gleichung}%        
			\renewcommand*{\itemautorefname}{Punkt}%
                }%  
            % Fix to getting autorefs for subfigures right (thanks to Belinda Vogt for changing the definition)
            \providecommand{\subfigureautorefname}{\figureautorefname}%             
    }{\relax}
\makeatother


% ****************************************************************************************************
% 7. Last calls before the bar closes
% ****************************************************************************************************

% ********************************************************************
% Development Stuff
% ********************************************************************
\listfiles
%\PassOptionsToPackage{l2tabu,orthodox,abort}{nag}
%   \usepackage{nag}
%\PassOptionsToPackage{warning, all}{onlyamsmath}
%   \usepackage{onlyamsmath}

% ********************************************************************
% Last, but not least...
% ********************************************************************
\usepackage{microtype} % Improves typography, reduces overflow via character protrusion/expansion
\usepackage[dottedtoc]{classicthesis} 
\usepackage{lmodern}

% ****************************************************************************************************

% ****************************************************************************************************
% 8. Further adjustments (experimental)
% ****************************************************************************************************
% ********************************************************************
% Changing the text area
% ********************************************************************
\linespread{1.05} % a bit more for Palatino
% \areaset[current]{312pt}{761pt} % 686 (factor 2.2) + 33 head + 42 head \the\footskip
% \setlength{\marginparwidth}{7em}%
% \setlength{\marginparsep}{2em}%

% ********************************************************************
% Fix text/equation overflow into margins
% ********************************************************************
\tolerance=1500                    % Allow looser spacing (default 200)
\emergencystretch=1.5em            % Extra stretch to avoid overflow
\hfuzz=1.5pt                       % Suppress warnings for minor overflow
\hbadness=1500   

% ********************************************************************
% Using different fonts
% ********************************************************************
%\usepackage[oldstylenums]{kpfonts} % oldstyle notextcomp
%\usepackage[osf]{libertine}
%\usepackage[light,condensed,math]{iwona}
%\renewcommand{\sfdefault}{iwona}
%\usepackage{lmodern} % <-- no osf support :-(
%\usepackage{cfr-lm} % 
%\usepackage[urw-garamond]{mathdesign} <-- no osf support :-(
%\usepackage[default,osfigures]{opensans} % scale=0.95 
%\usepackage[sfdefault]{FiraSans}
% ****************************************************************************************************

% ****************************************************************************************************
% X. My Custom Stuff
% ****************************************************************************************************
\usepackage{nomencl}
\usepackage{setspace}        % for \onehalfspacing
\usepackage{svg}
\usepackage{epigraph}
\usepackage{float}
\usepackage[section]{placeins} % keep floats within their sections/chapters
\usepackage{booktabs}
\usepackage{import}
\usepackage{array}
\usepackage{ragged2e}
\usepackage{longtable}
\usepackage{multirow}
\setlength{\LTleft}{0pt} % Fix longtable shift on page breaks in ClassicThesis
\setlength{\LTright}{\fill}
\usepackage{tikz}
\usepackage{lipsum}
% \usepackage[scaled=0.85]{beramono} % reload, after lmodern
% \renewcommand{\ttdefault}{fvm}     % make sure \ttfamily = Bera Mono
\usepackage[table]{xcolor}
\usepackage[most]{tcolorbox}
\usepackage[para]{threeparttable}
\usepackage{adjustbox}

\usepackage{xkeyval}     % for \define@key[FigTools]{...}, \setkeys
\usepackage{ifthen}      % for \ifthenelse and \equal
\usepackage{ifoddpage}   % for \checkoddpage and \ifoddpage

% Algorithms
\usepackage{algorithm}
\usepackage{algpseudocode}
% Each algorithm restarts its line numbers; include the float in PDF destinations.
\counterwithin*{ALG@line}{algorithm}

\usepackage{url}

% ********************************************************************
% CLEVEREF SETUP
% ********************************************************************
\usepackage[nameinlink,capitalize]{cleveref}
% Older LaTeX formats lose all but the first letter of acronym labels through
% cleveref's amsmath adapter. Use the complete key for acronym first-use links.
% https://github.com/oetiker/acronym/issues/62
\makeatletter
\IfFormatAtLeastTF{2025-06-01}{}{%
    \patchcmd{\AC@verridelabel}
        {\ifdefined\ltx@label\ltx@label}
        {\ifdefined\cref@label\cref@label}{}{}%
}
\makeatother

% 1. GENERIC SETTINGS
% -------------------
% Capitalize all references (even mid-sentence) to match Engineering standards
% e.g., "as seen in Fig. 1" rather than "as seen in fig. 1"

% 2. SPECIFIC DEFINITIONS
% -----------------------

% Equations: "Eq." and "Eqs."
\crefname{equation}{Eq.}{Eqs.}
\Crefname{equation}{Eq.}{Eqs.}
% Optional: Remove parentheses around eq number (e.g., "Eq. 1" instead of "Eq. (1)")
% If you prefer standard math format "Eq. (1)", comment out the line below.
\creflabelformat{equation}{#2#1#3}

% Figures: "Fig." and "Figs."
\crefname{figure}{Fig.}{Figs.}
\Crefname{figure}{Fig.}{Figs.}

% Tables: "Table" and "Tables" (Usually fully written out)
\crefname{table}{Table}{Tables}
\Crefname{table}{Table}{Tables}

% Code Listings (from listings package): "Listing" and "Listings"
\crefname{lstlisting}{Listing}{Listings}
\Crefname{lstlisting}{Listing}{Listings}

% Algorithms (from algorithm package): "Algorithm" and "Algorithms"
\crefname{algorithm}{Algorithm}{Algorithms}
\Crefname{algorithm}{Algorithm}{Algorithms}

% Sections/Chapters: "Section", "Chapter", "Part"
\crefname{part}{Part}{Parts}
\Crefname{part}{Part}{Parts}
\renewcommand{\thepart}{\Roman{part}}
% A ClassicThesis part ships its title page before \part returns.
\makeatletter
\newcommand{\labeledpart}[2]{%
    \GetTitleString{#1}\let\@currentlabelname\GetTitleStringResult
    \part[#1]{#1\label{#2}}%
}
\makeatother

\crefname{chapter}{Chapter}{Chapters}
\Crefname{chapter}{Chapter}{Chapters}

\crefname{section}{Section}{Sections}
\Crefname{section}{Section}{Sections}

\crefname{subsection}{Section}{Sections} % Map subsections to "Section"
\Crefname{subsection}{Section}{Sections}

\crefname{subsubsection}{Section}{Sections} % Map sub-subsections to "Section"
\Crefname{subsubsection}{Section}{Sections}

% Appendices
\crefname{appendix}{Appendix}{Appendices}
\Crefname{appendix}{Appendix}{Appendices}

% Enumerated Lists (items)
\crefname{enumi}{Item}{Items}
\Crefname{enumi}{Item}{Items}
\crefname{enumii}{Item}{Items}
\Crefname{enumii}{Item}{Items}
\crefname{enumiii}{Item}{Items}
\Crefname{enumiii}{Item}{Items}

% url penalties
\setcounter{biburllcpenalty}{7000}
\setcounter{biburlucpenalty}{8000}

% your column types
\newcolumntype{L}[1]{>{\RaggedRight\arraybackslash}p{#1}}
\newcolumntype{R}[1]{>{\RaggedLeft\arraybackslash}p{#1}}

% tcolorbox style
% Define color themes
\definecolor{lucasredlight}{RGB}{255,230,230}
\definecolor{lucasreddark}{RGB}{180,0,0}

\definecolor{lucasbluelight}{RGB}{225,235,255}
\definecolor{lucasbluedark}{RGB}{0,60,150}

% A key to choose theme = red or blue
\tcbset{
  lucasbox/.style={
    enhanced,
    fonttitle=\bfseries,
    title filled,
    boxrule=1pt,
    toptitle=2mm,
    bottomtitle=1mm
  },
  lucasboxtheme/red/.style={
    colback=lucasredlight,
    colframe=lucasreddark,
    coltitle=white,
    colbacktitle=lucasreddark
  },
  lucasboxtheme/blue/.style={
    colback=lucasbluelight,
    colframe=lucasbluedark,
    coltitle=white,
    colbacktitle=lucasbluedark
  }
}

% Default theme = red
\newenvironment{importantbox}[2][red]
  {\begin{tcolorbox}[lucasbox, lucasboxtheme/#1, title=#2]}
  {\end{tcolorbox}}

\DeclareMathOperator*{\argmax}{arg\,max}

\newcommand{\quotegraffito}[2]{\graffito{#1\\--~\emph{#2}}}

% grey midrule for trial block boundaries
\newcommand{\graymidrule}{%
  \arrayrulecolor{gray!60}%
  \midrule
  \arrayrulecolor{black}%
}

% Method labels (math and text safe)
\newcommand{\mBase}{\text{\textsc{Base}}}   % pretrained / SFT base policy
\newcommand{\mRL}{\text{\textsc{RL}}}       % pure RL from scratch
\newcommand{\mResRL}{\text{\textsc{ResRL}}} % residual RL on Base
